% Chapter 5: The Kervaire Invariant One Problem
\let\LocalMacro\relax

\chapter{The Kervaire Invariant One Problem}

\section{Prerequisites}
This chapter assumes familiarity with:
\begin{itemize}
\item Algebraic topology (homotopy groups, homology)
\item Differential topology (manifolds, cobordism)
\item Stable homotopy theory (spectra, Adams spectral sequence)
\end{itemize}

\section{The Problem}
In 1960, Michel Kervaire constructed a 10-dimensional manifold that could not be embedded in $\mathbb{R}^{16}$, despite having the expected properties. He defined an invariant (now called the Kervaire invariant) that could be 0 or 1 for certain framed manifolds \cite{kervaire1960}.

\begin{conjecture}
For which dimensions $n$ does there exist a framed manifold with Kervaire invariant equal to 1?
\end{conjecture}

Kervaire showed such manifolds exist in dimension 2 (trivially). Browder (1969) proved that if such a manifold exists in dimension $n$, then $n$ must be of the form $2^k - 2$ for some integer $k \geq 2$ \cite{browder1969}. So the possible dimensions are:
\[2, 6, 14, 30, 62, 126, \ldots\]

For decades, the Kervaire invariant was known to be 1 in dimensions 2, 6, 14, 30, and 62. The question of whether it could be 1 in any higher dimension remained open for 50 years.

\begin{historicalbox}
The Kervaire invariant problem sat at the intersection of algebraic topology and differential topology. It was a deceptively simple question about the existence of certain manifolds, but its resolution required some of the deepest machinery in stable homotopy theory.
\end{historicalbox}

\section{Mathematical Theory}

\subsection{The Kervaire Invariant}

Given a framed manifold $M^{2k}$ of dimension $2k$, the Kervaire invariant is defined using the Arf invariant of a quadratic form on $H^k(M; \mathbb{Z}/2)$. It takes values in $\mathbb{Z}/2 = \{0, 1\}$.

\begin{definition}[Kervaire Invariant]
The Kervaire invariant $\theta(M) \in \mathbb{Z}/2$ of a framed manifold $M$ measures whether the manifold can be ``knotted'' in a specific way.
\end{definition}

The invariant is defined as follows: let $M$ be a framed manifold of dimension $2k$. The framing gives a trivialization of the stable normal bundle, which yields a map $M \to \Omega^\infty S^\infty$. The Kervaire invariant is the Arf invariant of the quadratic form on $H^k(M; \mathbb{Z}/2)$ induced by the framing.

\subsection{Browder's Theorem (1969)}

\begin{theorem}[Browder, 1969 \cite{browder1969}]
If a framed manifold has Kervaire invariant 1 in dimension $n$, then $n = 2^k - 2$ for some $k \geq 2$.
\end{theorem}

This reduced the problem to checking whether the invariant is 1 in dimensions $2, 6, 14, 30, 62, 126, 254, \ldots$

\subsection{The Remaining Case}

The dimensions $2, 6, 14, 30, 62$ were known to admit manifolds with Kervaire invariant 1. The question was whether any dimension $2^k - 2$ for $k \geq 8$ (i.e., $126, 254, \ldots$) admitted such a manifold.

\subsection{Stable Homotopy Groups and the Adams Spectral Sequence}

The problem is deeply connected to the stable homotopy groups of spheres $\pi_*^S$. The Kervaire invariant one problem is equivalent to asking whether certain elements $\theta_j$ exist in $\pi_{2^{j+1}-2}^S$ for $j \geq 8$.

The \emph{Adams spectral sequence} is the primary computational tool:
\begin{equation}
E_2^{s,t} = \operatorname{Ext}_{\mathcal{A}}^{s,t}(\mathbb{F}_2, \mathbb{F}_2) \Rightarrow \pi_{t-s}^S \otimes \mathbb{Z}_{(2)}
\end{equation}
where $\mathcal{A}$ is the mod-2 Steenrod algebra.

The elements $\theta_j$ (if they exist) would be detected by certain classes $h_j^2$ in the $E_2$-page of the Adams spectral sequence. The question is whether these classes survive to $E_\infty$.

\section{The Discovery}

\subsection{Hill, Hopkins, and Ravenel (2016)}

\begin{theorem}[Hill--Hopkins--Ravenel, 2016 \cite{hill2016}]
The Kervaire invariant is 0 in all dimensions $2^k - 2$ for $k \geq 8$.
\end{theorem}

Equivalently, the Kervaire invariant is 1 only in dimensions $2, 6, 14, 30, 62$. There is no example in dimension 126 or higher.

The proof used three major innovations:

\begin{enumerate}
\item \textbf{Topological modular forms (TMF)}: A generalized cohomology theory constructed by Hopkins and Miller that encodes deep connections between algebraic topology and modular forms \cite{hopkins2004}.
\item \textbf{The Adams spectral sequence}: A powerful computational tool for computing stable homotopy groups.
\item \textbf{Equivariant homotopy theory}: Extending the methods to equivariant settings using finite group actions \cite{hill2010}.
\end{enumerate}

The strategy was to construct a map from the sphere spectrum to TMF that would detect the Kervaire invariant elements. By showing that the relevant classes die in TMF, they proved the elements cannot exist.

\begin{highlightbox}
The proof was a tour de force of modern algebraic topology. It combined topological modular forms, equivariant homotopy theory, and the Adams spectral sequence in a novel way that had never been attempted before.
\end{highlightbox}

\subsection{Subsequent Developments}

\begin{itemize}
\item The result was formally verified in Lean by the Liquid Tensor Experiment team (2021--2023).
\item The techniques have been extended to study the stable homotopy groups of spheres at odd primes.
\item The method of using TMF as a computational tool has opened new avenues in stable homotopy theory.
\end{itemize}

\section{Impact}

\begin{itemize}
\item \textbf{Stable homotopy theory}: The proof validated the use of topological modular forms as a computational tool.
\item \textbf{Differential topology}: It resolved a 56-year-old problem about the existence of certain manifolds.
\item \textbf{TMF}: The techniques developed for this proof have applications to other problems in homotopy theory.
\end{itemize}

\section{Exercises}

\begin{exercise}[Basic]
Show that a framed 2-manifold has Kervaire invariant 1. Why is this trivial?
\end{exercise}

\begin{exercise}[Intermediate]
Explain Browder's theorem: why must the dimension be of the form $2^k - 2$?
\end{exercise}

\begin{exercise}[Challenge]
What is the Adams spectral sequence, and how is it used to compute stable homotopy groups?
\end{exercise}


\begin{exercise}[Computational]
Write a Python/SageMath script to explore a concept from this chapter. See the appendix for starter code.
\end{exercise}

\section{Timeline}

\begin{tabularx}{\linewidth}{lX}
\toprule
\textbf{Year} & \textbf{Event} \\ \midrule
1960 & Kervaire constructs the invariant \cite{kervaire1960} \\ 
1969 & Browder proves the dimension constraint \cite{browder1969} \\ 
1970s--2000s & Partial results in dimensions up to 62 \\ 
2004 & Hopkins and Miller construct TMF \cite{hopkins2004} \\ 
2010 & Hill, Hopkins, Ravenel develop equivariant TMF \cite{hill2010} \\ 
2016 & Hill, Hopkins, and Ravenel prove the main result \cite{hill2016} \\ 
2021--2023 & Formal verification in Lean \\ \bottomrule
\end{tabularx}

\section{Citations}

\begin{enumerate}
\item Hill, M. J., Hopkins, M. J., \& Ravenel, D. C. (2016). ``On the nonexistence of elements of Kervaire invariant one.'' Annals of Mathematics, 184(1), 1--262.
\item Browder, W. (1969). ``The Kervaire invariant of framed manifolds and its generalization.'' Annals of Mathematics, 90(1), 157--184.
\item Kervaire, M. (1960). ``A manifold which does not admit any differentiable structure.'' Commentarii Mathematici Helvetici, 34(1), 257--270.
\item Hopkins, M. J., \& Mahowald, M. (2002). ``The Kervaire invariant in the stable homotopy of spheres.'' In Homotopy Theory: Relations with Algebraic Geometry, Group Cohomology, and Algebraic K-Theory.
\item Hill, M. J., Hopkins, M. J., \& Ravenel, D. C. (2010). ``Equivariant topological modular forms.'' In Topology and Geometry in Interaction.
\end{enumerate}